\begin{frame}{2.1절 정리}
  \begin{block}{핵심}
    \kb{ε-greedy}는 ``탐험의 양을 하나의 확률 $\varepsilon$ 으로
    조종한다''는, 이 책에서 첫 번째로 배우는 정책.
  \end{block}
  \begin{itemize}
    \item $\varepsilon=0$: ``한 번의 나쁜 운''에 최적해 영원히
      잃을 수 있음(2000스텝 후회 1754).
    \item $\varepsilon$ 지나치면: 이미 찾은 최선 팔을 외면한 채
      보상 낭비(후회 502).
    \item \textbf{탐험에도 비용이 있다} --- 그 비용과 ``탐험 부족''의
      대가를 줄이면 \textbf{가운데 어딘가}에 최적 $\varepsilon$.
  \end{itemize}
  \vskip0.5em
  이 ``탐험의 비용'' 감각은 2.2절(UCB·낙관적 초기화)에서
  \textbf{불확실성에 비례해} 탐험을 조절하는 방식으로 정교화되고,
  Week 5(몬테카를로)에서 \textbf{후회(regret)}로 정식 정의.
\end{frame}
